\section{Blocked-RFP construction for simple MPDOs}

\begin{definition}[Local simple-MPDO structure]
    \label{def:simple_mpdo_local_structure_data}
    \lean{MPOTensor.SimpleMPDOLocalStructureData}
    \leanok
    \uses{thm:sal_implies_eta_structure, thm:sal_zcl_implies_rank_one_t}
    This structure records the local Appendix~C.2 ingredients: a normalized
    three-site reduced state with equality in strong subadditivity, the
    resulting local $\eta$-structure, and the rank-one conclusion for the
    auxiliary primitive matrix $T$, supplied here through the conditional
    rank-one criterion of Theorem~\ref{thm:sal_zcl_implies_rank_one_t}.
\end{definition}

\begin{theorem}[Local simple-MPDO structure from SAL--ZCL hypotheses]
    \label{thm:simple_mpdo_local_structure_data_of_sal_zcl}
    \lean{MPOTensor.SimpleMPDOLocalStructureData.ofSALZCL}
    \lean{MPOTensor.SimpleMPDOLocalStructureData.ofSALZCLOfPosSemidef}
    \leanok
    \uses{def:simple_mpdo_local_structure_data, thm:sal_implies_eta_structure,
        thm:sal_zcl_implies_rank_one_t, thm:psd_trace_powers_rank_one_t}
    The hypotheses of Lemmas~C.2, C.4, and~C.5 determine the local simple-MPDO
    structure. There are two forms. The conditional form assumes the
    rank-one implication for $T$. The corrected form assumes instead that
    $T\ge 0$ as a matrix, with $\operatorname{tr}(T)=1$ and
    $\operatorname{tr}(T^m)=1$ for every $m>0$, and then applies
    Theorem~\ref{thm:psd_trace_powers_rank_one_t}.
\end{theorem}

\begin{definition}[Simple-MPDO blocked-RFP structure]
    \label{def:simple_mpdo_blocked_rfp_data}
    \lean{MPOTensor.SimpleMPDOBlockedRFPData}
    \leanok
    \uses{def:simple_mpdo_local_structure_data, def:mpdo_has_commuting_form,
        def:mpo_is_zcl}
    For an MPO tensor $K$, this records the local Appendix~C.2 ingredients, the
    commuting-form property, and zero correlation length. The local component
    consists of the hypotheses used in the entropy-side argument for commuting form.
\end{definition}

\begin{theorem}[Blocked-RFP structure from SAL--ZCL and commuting form]
    \label{thm:simple_mpdo_blocked_rfp_data_of_sal_zcl_and_commuting_form}
    \lean{MPOTensor.SimpleMPDOBlockedRFPData.ofSALZCLAndCommutingForm}
    \lean{MPOTensor.SimpleMPDOBlockedRFPData.ofSALZCLAndCommutingFormOfPosSemidef}
    \leanok
    \uses{def:simple_mpdo_blocked_rfp_data,
        thm:simple_mpdo_local_structure_data_of_sal_zcl}
    The SAL--ZCL hypotheses together with a commuting-form witness determine
    the simple-MPDO blocked-RFP structure. In the PSD-corrected form, the
    matrix input is $T\ge0$, $\operatorname{tr}(T)=1$, and
    $\operatorname{tr}(T^m)=1$ for all $m>0$, rather than an abstract
    Perron--Frobenius rank-one assumption for $T$.
\end{theorem}

\begin{theorem}[Blocked-RFP structure from $\eta$-local structure]
    \label{thm:simple_mpdo_blocked_rfp_data_of_eta_local_structure}
    \lean{MPOTensor.SimpleMPDOBlockedRFPData.ofEtaLocalStructure}
    \leanok
    \uses{def:simple_mpdo_local_structure_data,
        def:simple_mpdo_blocked_rfp_data,
        def:mpdo_eta_local_structure_data}
    Given the local Appendix~C.2 data, ZCL, and an $\eta$-local product
    $\sigma^{(N)}(K)=c_N\prod_n B_{n,n+1}$ with $c_N>0$,
    $B_{n,n+1}\ge0$, and $[B_{i,i+1},B_{j,j+1}]=0$, one obtains the
    simple-MPDO blocked-RFP structure for $K$.
\end{theorem}

\begin{proof}\leanok
    The $\eta$-local product gives the commuting-form property for $K$; the local
    Appendix~C.2 hypotheses and ZCL fill the other two fields.
\end{proof}

\begin{theorem}[Blocked-RFP structure from SAL--ZCL and $\eta$-local structure]
    \label{thm:simple_mpdo_blocked_rfp_data_of_sal_zcl_and_eta_local_structure}
    \lean{MPOTensor.SimpleMPDOBlockedRFPData.ofSALZCLAndEtaLocalStructure}
    \lean{MPOTensor.SimpleMPDOBlockedRFPData.ofSALZCLAndEtaLocalStructureOfPosSemidef}
    \leanok
    \uses{thm:simple_mpdo_local_structure_data_of_sal_zcl,
        thm:simple_mpdo_blocked_rfp_data_of_eta_local_structure}
    Suppose Lemma~C.2 gives the Markov decomposition, Lemma~C.4 gives the
    local $\eta$-structure and a primitive matrix $T$ with
    $\operatorname{tr}(T)=1$, and ZCL together with Lemma~C.5 gives
    $\operatorname{tr}(T^m)$ independent of $m$ and the conditional rank-one
    criterion for $T$. Suppose also that the assembled $\eta$-local structure gives
    $\sigma^{(N)}(K)=c_N\prod_{n=1}^N B_{n,n+1}$ with $c_N>0$,
    $B_{n,n+1}\ge 0$, and $[B_{i,i+1},B_{j,j+1}]=0$ for all $i,j$.
    Then these hypotheses determine the simple-MPDO blocked-RFP structure for $K$.
    In the PSD-corrected form, the conditional rank-one criterion is replaced
    by $T\ge0$, $\operatorname{tr}(T)=1$, and
    $\operatorname{tr}(T^m)=1$ for all $m>0$.
\end{theorem}

\begin{proof}\leanok
    First form the local structure from the Lemma~C.2 Markov decomposition,
    the Lemma~C.4 local $\eta$-structure and primitivity statement, and the
    Lemma~C.5 conditional rank-one criterion for $T$. The $\eta$-local product
    formula supplies the commuting-form component, and ZCL supplies the
    remaining component of the blocked-RFP structure.
\end{proof}

\begin{theorem}[Zero correlation length gives a blocked fusion-isometry witness]
    \label{thm:structural_implies_rfp_blocked}
    \lean{MPOTensor.structural_implies_rfp_blocked}
    \leanok
    \uses{def:mpo_is_zcl, def:mpdo_fusion_isometry, def:mpdo_rfp_via_fusion,
        thm:mpdo_rfp_via_fusion_iff}
    If an MPO tensor $K$ has zero correlation length, then it has a blocked
    fusion-isometry witness at size $2$ in the sense of
    Definition~\ref{def:mpdo_fusion_isometry}.
\end{theorem}

\begin{proof}\leanok
    Zero correlation length is transfer-map idempotence.
    Applying the converse direction of
    Theorem~\ref{thm:mpdo_rfp_via_fusion_iff} to that RFP witness yields
    fusion-isometry data at every positive blocked size; evaluating it at size $2$
    produces the blocked witness.
\end{proof}

\begin{theorem}[Blocked-RFP data give a blocked fusion-isometry witness]
    \label{thm:structural_implies_rfp_blocked_of_data}
    \lean{MPOTensor.structural_implies_rfp_blocked_of_data}
    \leanok
    \uses{def:simple_mpdo_blocked_rfp_data, thm:structural_implies_rfp_blocked}
    Given the blocked-RFP data for an MPO tensor $K$, one obtains a blocked
    fusion-isometry witness at size $2$.
\end{theorem}

\begin{proof}\leanok
    Apply the zero-correlation-length theorem to the ZCL component of the
    blocked-RFP data.
\end{proof}

\begin{lemma}[Blocked-RFP data give the fusion formulation]
    \label{lem:simple_mpdo_blocked_rfp_data_rfp_via_fusion}
    \lean{MPOTensor.SimpleMPDOBlockedRFPData.isRFP_via_fusion}
    \leanok
    \uses{def:simple_mpdo_blocked_rfp_data, def:mpdo_rfp_via_fusion}
    Given the blocked-RFP data for an MPO tensor $K$, one obtains the
    transfer-map fusion formulation of idempotence.
\end{lemma}

\begin{proof}\leanok
    The implication used here is
    \[
        \operatorname{ZCL}(K)
        \Longrightarrow \operatorname{RFP}(K)
        \Longrightarrow \operatorname{RFP}_{\mathrm{fusion}}(K).
    \]
\end{proof}

\begin{theorem}[Simple-MPDO RFP chain from blocked-RFP data]
    \label{thm:simple_mpdo_rfp_chain_of_data}
    \lean{MPOTensor.simple_mpdo_rfp_chain_of_data}
    \leanok
    \uses{def:simple_mpdo_blocked_rfp_data,
        thm:structural_implies_rfp_blocked,
        lem:simple_mpdo_blocked_rfp_data_rfp_via_fusion}
    Given the blocked-RFP data for an MPO tensor $K$, one obtains the
    GSNNCH-with-ZCL case, the blocked fusion-isometry witness at size $2$, and
    the transfer-map fusion formulation of idempotence.
\end{theorem}

\begin{proof}\leanok
    The commuting-form and zero-correlation-length assumptions give the
    GSNNCH-with-ZCL case. Zero correlation length also gives the blocked
    fusion-isometry witness and the transfer-map fusion formulation.
\end{proof}

\begin{lemma}[Size-$2$ fusion witness from $\eta$-local structure]
    \label{lem:eta_local_structure_fusion_isometry_two}
    \lean{MPOTensor.EtaLocalStructureData.fusion_isometry_data_two}
    \leanok
    \uses{def:mpdo_eta_local_structure_data, def:mpo_is_zcl,
        thm:structural_implies_rfp_blocked}
    If the $\eta$-local structure has been constructed for an MPO tensor $K$,
    then zero correlation length gives the blocked fusion-isometry witness at
    size $2$.
\end{lemma}

\begin{proof}\leanok
    The implication used here is
    \[
        \operatorname{ZCL}(K)
        \Longrightarrow \operatorname{RFP}(K)
        \Longrightarrow \operatorname{FusionWitness}_{2}(K).
    \]
\end{proof}

\begin{lemma}[Fusion formulation from $\eta$-local structure]
    \label{lem:eta_local_structure_rfp_via_fusion}
    \lean{MPOTensor.EtaLocalStructureData.isRFP_via_fusion}
    \leanok
    \uses{def:mpdo_eta_local_structure_data, def:mpo_is_zcl,
        def:mpdo_rfp_via_fusion}
    If the $\eta$-local structure has been constructed for an MPO tensor $K$,
    then zero correlation length gives the transfer-map fusion formulation of
    idempotence.
\end{lemma}

\begin{proof}\leanok
    This is the implication
    \[
        \operatorname{ZCL}(K)
        \Longrightarrow \operatorname{RFP}(K)
        \Longrightarrow \operatorname{RFP}_{\mathrm{fusion}}(K).
    \]
\end{proof}

\begin{theorem}[Simple-MPDO RFP chain from $\eta$-local structure]
    \label{thm:simple_mpdo_rfp_chain_of_eta_local_structure}
    \lean{MPOTensor.EtaLocalStructureData.simple_mpdo_rfp_chain}
    \lean{MPOTensor.simple_mpdo_rfp_chain_of_etaLocalStructure}
    \leanok
    \uses{def:mpdo_eta_local_structure_data, def:mpo_is_zcl,
        thm:mpdo_eta_local_structure_is_gsnnch_zcl,
        lem:eta_local_structure_fusion_isometry_two,
        lem:eta_local_structure_rfp_via_fusion}
    If the $\eta$-local structure has been constructed for an
    MPO tensor $K$, then zero correlation length gives the GSNNCH-with-ZCL
    case, the blocked fusion-isometry witness at size $2$, and the
    transfer-map fusion formulation of idempotence.
\end{theorem}

\begin{proof}\leanok
    The $\eta$-local structure supplies the commuting nearest-neighbor product
    form. The remaining two implications are
    \[
        \operatorname{ZCL}(K) \Longrightarrow
        \operatorname{RFP}(K) \Longrightarrow
        \operatorname{RFP}_{\mathrm{fusion}}(K)
        \quad\text{and}\quad
        \operatorname{ZCL}(K) \Longrightarrow
        \operatorname{RFP}(K) \Longrightarrow
        \operatorname{FusionWitness}_{2}(K).
    \]
\end{proof}

\begin{theorem}[Simple-MPDO RFP chain from SAL--ZCL and $\eta$-local structure]
    \label{thm:simple_mpdo_rfp_chain_of_sal_zcl_and_eta_local_structure}
    \lean{MPOTensor.simple_mpdo_rfp_chain_of_sal_zcl_and_etaLocalStructure}
    \lean{MPOTensor.simple_mpdo_rfp_chain_of_sal_zcl_and_etaLocalStructure_of_posSemidef}
    \leanok
    \uses{thm:simple_mpdo_blocked_rfp_data_of_sal_zcl_and_eta_local_structure,
        thm:simple_mpdo_rfp_chain_of_data}
    Assume the local SAL--ZCL hypotheses of Lemmas~C.2, C.4, and~C.5 and the assembled
    $\eta$-local product form, with the conditional rank-one criterion for $T$,
    $\sigma^{(N)}(K)=c_N\prod_{n=1}^N B_{n,n+1}$ with $c_N>0$,
    $B_{n,n+1}\ge 0$, and $[B_{i,i+1},B_{j,j+1}]=0$. Then
    $K$ is GSNNCH with ZCL, admits a blocked fusion-isometry witness at size
    $2$, and satisfies the transfer-map fusion formulation of idempotence. The
    PSD-corrected form uses $T\ge0$ together with
    $\operatorname{tr}(T^m)=1$ for all $m>0$ in place of the conditional
    matrix rank-one hypothesis.
\end{theorem}

\begin{proof}\leanok
    Combine the local SAL--ZCL hypotheses and the $\eta$-local product form into the
    blocked-RFP structure, then apply
    Theorem~\ref{thm:simple_mpdo_rfp_chain_of_data}.
\end{proof}

\begin{theorem}[Simple-MPDO RFP chain from commuting form and ZCL]
    \label{thm:simple_mpdo_rfp_chain}
    \lean{MPOTensor.simple_mpdo_rfp_chain}
    \leanok
    \uses{def:mpdo_has_commuting_form, def:mpo_is_zcl,
        thm:mpdo_is_gsnnch_zcl_iff_has_commuting_form_and_is_zcl,
        thm:structural_implies_rfp_blocked, def:mpdo_rfp_via_fusion,
        thm:mpdo_rfp_via_fusion_iff}
    From the commuting-form property and zero correlation length one
    simultaneously recovers the GSNNCH-with-ZCL criterion, the
    blocked fusion-isometry witness at size $2$, and the transfer-map fusion
    formulation of idempotence.
\end{theorem}

\begin{proof}\leanok
    The commuting-form and zero-correlation-length hypotheses give the GSNNCH--ZCL
    branch via Theorem~\ref{thm:mpdo_is_gsnnch_zcl_iff_has_commuting_form_and_is_zcl}.
    The blocked witness is Theorem~\ref{thm:structural_implies_rfp_blocked}. The
    full transfer-map fusion formulation follows by applying the converse direction
    of Theorem~\ref{thm:mpdo_rfp_via_fusion_iff} to zero correlation length, which
    is transfer-map idempotence.
\end{proof}

\begin{remark}[SAL--ZCL implication for the GSNNCH--ZCL criterion]
    \label{rem:simple_mpdo_rfp_chain_gap}
    \notready
    \uses{def:simple_mpdo_local_structure_data,
        def:simple_mpdo_blocked_rfp_data, thm:simple_mpdo_rfp_chain}
    The implication from SAL and ZCL alone to the conclusion of Appendix~C.2
    requires a derivation of the commuting-form property from the local
    simple-MPDO structure of Lemmas~C.2, C.4, and~C.5. Once that derivation is available,
    Theorem~\ref{thm:simple_mpdo_rfp_chain} gives the GSNNCH--ZCL branch and
    the blocked fusion-isometry witness. This missing derivation is documented
    in \path{docs/paper-gaps/cpgsv17_mpdo_sal_zcl_eta_local_structure.tex}
    (issue~\#823).
\end{remark}

